% experiments.tex
\section{Experiments}
\label{sec:experiments}

We evaluate KnowFeat on both public tabular benchmarks and a real-world anti-money laundering (AML) dataset. Our experiments aim to answer four research questions:
\begin{itemize}
    \item \textbf{RQ1}: Does KnowFeat improve predictive performance on public datasets?
    \item \textbf{RQ2}: How does KnowFeat perform on a real-world AML task with severe class imbalance?
    \item \textbf{RQ3}: What is the contribution of domain knowledge? (Ablation study)
    \item \textbf{RQ4}: How does KnowFeat balance feature quantity vs.\ quality?
\end{itemize}

\subsection{Experimental Setup}
\label{subsec:setup}

\subsubsection{Datasets}
We use eight datasets summarized in Table~\ref{tab:datasets}: seven public benchmarks from OpenML and one real-world AML dataset.

\begin{table}[t]
\centering
\caption{Dataset statistics.}
\label{tab:datasets}
\small
\renewcommand{\arraystretch}{1.2}
\setlength{\tabcolsep}{3pt}
\begin{tabular}{l r r r l}
\hline
\textbf{Dataset} & \textbf{Samples} & \textbf{Feat.} & \textbf{Pos\%} & \textbf{Domain} \\
\hline
credit\_g          & 1,000   & 20  & 70.0 & Finance \\
diabetes           & 768     & 8   & 34.9 & Healthcare \\
blood\_transfusion & 748     & 4   & 23.8 & Healthcare \\
heart              & 918     & 11  & 55.3 & Healthcare \\
bank\_marketing    & 45,211  & 16  & 11.7 & Marketing \\
breast\_w          & 683     & 9   & 35.0 & Healthcare \\
adult              & 48,842  & 14  & 23.9 & Census \\
SimECNY (AML)      & 200,000 & 93  & 2.24 & Finance \\
\hline
\end{tabular}
\end{table}

The seven public datasets are standard tabular benchmarks covering diverse domains. SimECNY is a real-world digital currency transaction dataset containing 200,000 transactions with 93 baseline features and 4,484 positive (suspicious) samples, representing a realistic AML scenario with severe class imbalance (2.24\% positive rate).

\subsubsection{Baselines}
We compare against five baselines:
\begin{itemize}
    \item \textbf{No-FE}: XGBoost on original features.
    \item \textbf{OpenFE}~\cite{zhang2024openfe}: Transformation-based AutoFE with LightGBM-based selection.
    \item \textbf{AutoFeat}~\cite{horn2020set}: Polynomial feature generation with L1 selection.
    \item \textbf{CAAFE}~\cite{hollmann2024llmsgeneratefinanciallyrelevantfeatures}: LLM-based iterative feature generation.
    \item \textbf{FeatLLM}~\cite{han2024featllm}: LLM-based rule extraction and binary feature generation.
\end{itemize}

\subsubsection{Evaluation Protocol}
All methods are evaluated using XGBoost (n\_estimators=200, max\_depth=6, learning\_rate=0.1) with stratified 80/20 train-test splits across 5 random seeds. For public datasets, the primary metric is AUC. For SimECNY, we report Recall, Precision, F1, and AUC with scale\_pos\_weight=auto ($\approx$43.6) to handle class imbalance. We additionally evaluate with scale\_pos\_weight=1 (no class balancing) to isolate the intrinsic quality of generated features.

\subsection{RQ1: Public Dataset Performance}
\label{subsec:rq1}

Table~\ref{tab:main_results} presents the AUC comparison on seven public datasets. KnowFeat achieves the best performance on 4 out of 7 datasets (credit\_g, blood\_transfusion, breast\_w, and adult) and ranks second on 2 datasets (diabetes and heart).

\begin{table*}[t]
\centering
\caption{AUC comparison on public datasets (mean$\pm$std over 5 seeds). \textbf{Bold}: best; \underline{underline}: second best. T/O = timeout.}
\label{tab:main_results}
\small
\renewcommand{\arraystretch}{1.2}
\setlength{\tabcolsep}{6pt}
\begin{tabular}{l c c c c c c}
\hline
\textbf{Dataset} & \textbf{No-FE} & \textbf{OpenFE} & \textbf{AutoFeat} & \textbf{CAAFE} & \textbf{FeatLLM} & \textbf{KnowFeat} \\
\hline
credit\_g          & .7670{\scriptsize$\pm$.021} & \underline{.7833{\scriptsize$\pm$.031}} & .7540{\scriptsize$\pm$.017} & .7596{\scriptsize$\pm$.023} & .7604{\scriptsize$\pm$.022} & \textbf{.7854{\scriptsize$\pm$.014}} \\
diabetes           & .8027{\scriptsize$\pm$.033} & \textbf{.8252{\scriptsize$\pm$.034}} & .8089{\scriptsize$\pm$.027} & .8032{\scriptsize$\pm$.030} & .8096{\scriptsize$\pm$.032} & \underline{.8190{\scriptsize$\pm$.032}} \\
blood\_transfusion & .6824{\scriptsize$\pm$.043} & .6654{\scriptsize$\pm$.040} & .6824{\scriptsize$\pm$.043} & .6779{\scriptsize$\pm$.043} & .6805{\scriptsize$\pm$.044} & \textbf{.7007{\scriptsize$\pm$.036}} \\
heart              & .9320{\scriptsize$\pm$.007} & \textbf{.9417{\scriptsize$\pm$.011}} & .9340{\scriptsize$\pm$.003} & .9302{\scriptsize$\pm$.006} & .9297{\scriptsize$\pm$.004} & \underline{.9379{\scriptsize$\pm$.006}} \\
bank\_marketing    & .9330{\scriptsize$\pm$.002} & \textbf{.9374{\scriptsize$\pm$.002}} & T/O & .9334{\scriptsize$\pm$.002} & .9337{\scriptsize$\pm$.002} & \underline{.9344{\scriptsize$\pm$.002}} \\
breast\_w          & .9926{\scriptsize$\pm$.005} & \underline{.9944{\scriptsize$\pm$.005}} & .9936{\scriptsize$\pm$.004} & .9924{\scriptsize$\pm$.006} & .9924{\scriptsize$\pm$.005} & \textbf{.9957{\scriptsize$\pm$.003}} \\
adult              & .9288{\scriptsize$\pm$.004} & \underline{.9300{\scriptsize$\pm$.003}} & .9254{\scriptsize$\pm$.004} & .9286{\scriptsize$\pm$.003} & .9289{\scriptsize$\pm$.003} & \textbf{.9301{\scriptsize$\pm$.003}} \\
\hline
\textbf{Avg.\ Rank} & 4.1 & \underline{2.1} & 4.4 & 4.9 & 4.1 & \textbf{1.4} \\
\hline
\end{tabular}
\end{table*}

KnowFeat achieves the best average rank (1.4), outperforming OpenFE (2.1). The largest improvements occur on datasets where domain knowledge provides strong semantic guidance: credit\_g (+0.018 over No-FE) and blood\_transfusion (+0.018 over No-FE).

OpenFE achieves the best results on 3 datasets (diabetes, heart, bank\_marketing) through exhaustive operator-operand search. However, OpenFE generates 34--50 features per dataset, while KnowFeat achieves competitive or superior performance with 12--46 features (see Section~\ref{subsec:rq4}).

CAAFE and FeatLLM perform close to the No-FE baseline on most datasets, suggesting that features generated without rigorous verification often introduce noise rather than signal.

\subsection{RQ2: AML Dataset Performance}
\label{subsec:rq2}

Table~\ref{tab:simecny} presents results on the SimECNY AML dataset with scale\_pos\_weight=auto. KnowFeat adds 18 knowledge-guided features to the 93-feature baseline, achieving the best results across all four metrics.

\begin{table}[t]
\centering
\caption{SimECNY AML results (spw=auto, 5 seeds).}
\label{tab:simecny}
\small
\renewcommand{\arraystretch}{1.2}
\setlength{\tabcolsep}{2pt}
\begin{tabular}{l r c c c c}
\hline
\textbf{Method} & $\boldsymbol{n_f}$ & \textbf{Recall} & \textbf{Prec.} & \textbf{F1} & \textbf{AUC} \\
\hline
No-FE     & 93   & .9663{\scriptsize$\pm$.010} & .9764{\scriptsize$\pm$.004} & .9713{\scriptsize$\pm$.005} & .9996 \\
OpenFE    & 123  & .9645{\scriptsize$\pm$.008} & .9752{\scriptsize$\pm$.004} & .9698{\scriptsize$\pm$.003} & .9996 \\
CAAFE     & 94   & .9681{\scriptsize$\pm$.008} & .9725{\scriptsize$\pm$.006} & .9703{\scriptsize$\pm$.005} & .9996 \\
FeatLLM   & 1124 & .9625{\scriptsize$\pm$.010} & .9728{\scriptsize$\pm$.005} & .9676{\scriptsize$\pm$.006} & .9995 \\
\textbf{Ours} & \textbf{111} & \textbf{.9871}{\scriptsize$\pm$.006} & \textbf{.9855}{\scriptsize$\pm$.001} & \textbf{.9863}{\scriptsize$\pm$.003} & \textbf{.9999} \\
\hline
\end{tabular}
\end{table}

KnowFeat's 18 domain-knowledge features yield substantial improvements: Recall improves by +2.1pp over No-FE (0.9871 vs.\ 0.9663), F1 by +1.5pp (0.9863 vs.\ 0.9713), and AUC reaches 0.9999. In contrast, OpenFE adds 30 features yet \emph{decreases} Recall, while FeatLLM adds 1,031 features and achieves the \emph{worst} performance across all metrics. This demonstrates that knowledge-guided features provide qualitatively different signal from syntactic transformations.

\subsubsection{Robustness Analysis (spw=1)}

To verify that KnowFeat's gains stem from intrinsic feature quality rather than interaction with class balancing, we re-evaluate all methods with scale\_pos\_weight=1 (Table~\ref{tab:simecny_spw1}).

\begin{table}[t]
\centering
\caption{SimECNY AML results (spw=1, 5 seeds).}
\label{tab:simecny_spw1}
\small
\renewcommand{\arraystretch}{1.2}
\setlength{\tabcolsep}{2pt}
\begin{tabular}{l r c c c c}
\hline
\textbf{Method} & $\boldsymbol{n_f}$ & \textbf{Recall} & \textbf{Prec.} & \textbf{F1} & \textbf{AUC} \\
\hline
No-FE     & 93   & .9409{\scriptsize$\pm$.011} & .9934{\scriptsize$\pm$.002} & .9664{\scriptsize$\pm$.006} & .9996 \\
OpenFE    & 123  & .9342{\scriptsize$\pm$.012} & .9912{\scriptsize$\pm$.002} & .9619{\scriptsize$\pm$.007} & .9995 \\
CAAFE     & 94   & .9382{\scriptsize$\pm$.012} & .9922{\scriptsize$\pm$.002} & .9644{\scriptsize$\pm$.007} & .9996 \\
FeatLLM   & 1124 & .9298{\scriptsize$\pm$.014} & .9912{\scriptsize$\pm$.002} & .9595{\scriptsize$\pm$.008} & .9995 \\
\textbf{Ours} & \textbf{111} & \textbf{.9683}{\scriptsize$\pm$.010} & \textbf{.9945}{\scriptsize$\pm$.004} & \textbf{.9812}{\scriptsize$\pm$.005} & \textbf{.9999} \\
\hline
\end{tabular}
\end{table}

Without class balancing, KnowFeat remains the clear winner, with Recall 0.9683 (+2.7pp over No-FE) and F1 0.9812 (+1.5pp). Importantly, KnowFeat is the \emph{only} method that outperforms the No-FE baseline---all other feature engineering methods \emph{degrade} Recall relative to the 93-feature baseline. Moreover, KnowFeat suffers the smallest Recall drop when switching from spw=auto to spw=1 ($\Delta = -$1.9pp), while FeatLLM drops the most ($\Delta = -$3.3pp). This confirms that KnowFeat's 18 features encode genuine discriminative signal that does not depend on class reweighting.

\subsection{RQ3: Ablation Study}
\label{subsec:rq3}

To quantify the contribution of domain knowledge, we compare the \textit{schema\_only} variant (Level 0: only dataset metadata) against the \textit{full} variant (all knowledge levels). Results on public datasets are shown in Table~\ref{tab:ablation}.

\begin{table}[t]
\centering
\caption{Ablation: schema\_only vs.\ full KnowFeat (AUC).}
\label{tab:ablation}
\small
\renewcommand{\arraystretch}{1.2}
\setlength{\tabcolsep}{3pt}
\begin{tabular}{l c c c}
\hline
\textbf{Dataset} & \textbf{schema\_only} & \textbf{full} & $\boldsymbol{\Delta}$ \\
\hline
credit\_g     & .7750{\scriptsize$\pm$.012} & \textbf{.7854}{\scriptsize$\pm$.014} & +.010 \\
blood\_trans.  & .6891{\scriptsize$\pm$.048} & \textbf{.7007}{\scriptsize$\pm$.036} & +.012 \\
heart         & .9304{\scriptsize$\pm$.005} & \textbf{.9379}{\scriptsize$\pm$.006} & +.008 \\
bank\_mkt.    & .9314{\scriptsize$\pm$.002} & \textbf{.9344}{\scriptsize$\pm$.002} & +.003 \\
adult         & .9292{\scriptsize$\pm$.003} & \textbf{.9301}{\scriptsize$\pm$.003} & +.001 \\
breast\_w     & \textbf{.9964}{\scriptsize$\pm$.002} & .9957{\scriptsize$\pm$.003} & $-$.001 \\
diabetes      & --- (0 feat.) & .8190{\scriptsize$\pm$.032} & --- \\
\hline
\end{tabular}
\end{table}

\textbf{Knowledge drives feature generation.} On the diabetes dataset, the schema-only variant generates \emph{zero} features that pass L3 verification---the LLM cannot create useful features from column names alone. With full knowledge, KnowFeat generates 24 accepted features and improves AUC by +0.016 over No-FE.

\textbf{Knowledge improves feature quality.} On five of the six remaining datasets, the full variant outperforms schema\_only, with gains ranging from +0.001 (adult) to +0.012 (blood\_transfusion). On blood\_transfusion, both variants generate 12 features, indicating that the improvement comes from feature \emph{quality}---domain knowledge guides the agent toward more predictive transformations.

\textbf{breast\_w exception.} On breast\_w, schema\_only slightly outperforms full ($-$0.001 AUC). We attribute this to the near-ceiling baseline (AUC $>$ 0.99), where additional features may introduce marginal noise.

\subsection{RQ4: Feature Quantity vs.\ Quality}
\label{subsec:rq4}

Table~\ref{tab:feature_count} compares the number of features generated by each method.

\begin{table}[t]
\centering
\caption{Number of features (incl.\ original).}
\label{tab:feature_count}
\small
\renewcommand{\arraystretch}{1.2}
\setlength{\tabcolsep}{2.5pt}
\begin{tabular}{l r r r r r r}
\hline
\textbf{Dataset} & \textbf{No-FE} & \textbf{OpenFE} & \textbf{Auto.} & \textbf{CAAFE} & \textbf{Feat.} & \textbf{Ours} \\
\hline
credit\_g     & 20  & 50  & 23  & 24  & 280  & 46  \\
diabetes      & 8   & 38  & 13  & 9   & 152  & 24  \\
blood\_trans. & 4   & 34  & 4   & 7   & 120  & 12  \\
heart         & 11  & 41  & 21  & 14  & 98   & 31  \\
bank\_mkt.    & 16  & 46  & T/O & 21  & 132  & 36  \\
breast\_w     & 9   & 39  & 20  & 10  & 107  & 16  \\
adult         & 14  & 44  & 165 & 19  & 170  & 34  \\
SimECNY       & 93  & 123 & --- & 94  & 1124 & 111 \\
\hline
\end{tabular}
\end{table}

FeatLLM generates by far the most features (98--1,124) yet achieves the worst or near-worst performance on most datasets. This confirms that \emph{more features $\neq$ better performance} when verification is absent. AutoFeat generates 165 features on adult yet ranks last among all methods (AUC 0.9254), further illustrating this point.

KnowFeat strikes an effective balance: it generates a moderate number of features (12--46 on public datasets, +18 on SimECNY) that have all passed the L1$\to$L2$\to$L3 pipeline. On SimECNY, the contrast is particularly stark: KnowFeat's 18 verified features improve F1 by +1.5pp, while FeatLLM's 1,031 unverified features \emph{decrease} F1 by $-$0.4pp. This demonstrates the effectiveness of quality-over-quantity feature engineering.

\subsection{Provenance and Interpretability}
\label{subsec:provenance}

Every KnowFeat-generated feature includes a provenance card recording: (1) the knowledge assets that inspired it; (2) a reasoning chain explaining \emph{why} the feature is relevant; and (3) the computation steps from raw fields to the final value. This traceability is essential for deployment in regulated industries. We illustrate this capability through detailed case studies in Section~\ref{sec:case_study}.
